
\documentclass{article}
\usepackage{colortbl}
\usepackage{makecell}
\usepackage{multirow}
\usepackage{supertabular}

\begin{document}

\newcounter{utterance}

\centering \large Interaction Transcript for game `hot\_air\_balloon', experiment `air\_balloon\_survival\_de\_negotiation\_hard', episode 0 with qwen3{-}next{-}80b{-}a3b{-}thinking{-}t1.0.
\vspace{24pt}

{ \footnotesize  \setcounter{utterance}{1}
\setlength{\tabcolsep}{0pt}
\begin{supertabular}{c@{$\;$}|p{.15\linewidth}@{}p{.15\linewidth}p{.15\linewidth}p{.15\linewidth}p{.15\linewidth}p{.15\linewidth}}
    \# & $\;$A & \multicolumn{4}{c}{Game Master} & $\;\:$B\\
    \hline

    \theutterance \stepcounter{utterance}  
    & & \multicolumn{4}{p{0.6\linewidth}}{
        \cellcolor[rgb]{0.9,0.9,0.9}{
            \makecell[{{p{\linewidth}}}]{
                \texttt{\tiny{[P1$\langle$GM]}}
                \texttt{Du nimmst an einem kollaborativen Verhandlungspiel Teil.} \\
\\ 
\texttt{Zusammen mit einem anderen Teilnehmer musst du dich auf eine Reihe von Gegenständen entscheiden, die behalten werden. Jeder von euch hat eine persönliche Verteilung über die Wichtigkeit der einzelnen Gegenstände. Jeder von euch hat eine eigene Meinung darüber, wie wichtig jeder einzelne Gegenstand ist (Gegenstandswichtigkeit). Du kennst die Wichtigkeitsverteilung des anderen Spielers nicht. Zusätzlich siehst du, wie viel Aufwand jeder Gegenstand verursacht.  } \\
\texttt{Ihr dürft euch nur auf eine Reihe von Gegenständen einigen, wenn der Gesamtaufwand der ausgewählten Gegenstände den Maximalaufwand nicht überschreitet:} \\
\\ 
\texttt{Maximalaufwand = 3624} \\
\\ 
\texttt{Hier sind die einzelnen Aufwände der Gegenstände:} \\
\\ 
\texttt{Aufwand der Gegenstände = \{"C33": 210, "A89": 291, "A48": 831, "A04": 111, "A41": 926, "C64": 827, "C49": 25, "A61": 121, "C04": 583, "B53": 766, "B85": 14, "B61": 559, "A45": 304, "A03": 989, "B21": 691\}} \\
\\ 
\texttt{Hier ist deine persönliche Verteilung der Wichtigkeit der einzelnen Gegenstände:} \\
\\ 
\texttt{Werte der Gegenstandswichtigkeit = \{"C33": 138, "A89": 583, "A48": 868, "A04": 822, "A41": 783, "C64": 65, "C49": 262, "A61": 121, "C04": 508, "B53": 780, "B85": 461, "B61": 484, "A45": 668, "A03": 389, "B21": 808\}} \\
\\ 
\texttt{Ziel:} \\
\\ 
\texttt{Dein Ziel ist es, eine Reihe von Gegenständen auszuhandeln, die dir möglichst viel bringt (d. h. Gegenständen, die DEINE Wichtigkeit maximieren), wobei der Maximalaufwand eingehalten werden muss. Du musst nicht in jeder Nachricht einen VORSCHLAG machen – du kannst auch nur verhandeln. Alle Taktiken sind erlaubt!} \\
\\ 
\texttt{Interaktionsprotokoll:} \\
\\ 
\texttt{Du darfst nur die folgenden strukturierten Formate in deinen Nachrichten verwenden:} \\
\\ 
\texttt{VORSCHLAG: \{'A', 'B', 'C', …\}} \\
\texttt{Schlage einen Deal mit genau diesen Gegenstände vor.} \\
\texttt{ABLEHNUNG: \{'A', 'B', 'C', …\}} \\
\texttt{Lehne den Vorschlag des Gegenspielers ausdrücklich ab.} \\
\texttt{ARGUMENT: \{'...'\}} \\
\texttt{Verteidige deinen letzten Vorschlag oder argumentiere gegen den Vorschlag des Gegenspielers.} \\
\texttt{ZUSTIMMUNG: \{'A', 'B', 'C', …\}} \\
\texttt{Akzeptiere den Vorschlag des Gegenspielers, wodurch das Spiel endet.} \\
\texttt{STRATEGISCHE ÜBERLEGUNGEN: \{'...'\}} \\
\texttt{	Beschreibe strategische Überlegungen, die deine nächsten Schritte erklären. Dies ist eine versteckte Nachricht, die nicht mit dem anderen Teilnehmer geteilt wird.} \\
\\ 
\texttt{Regeln:} \\
\\ 
\texttt{Du darst nur einen Vorschlag mit ZUSTIMMUNG akzeptieren, der vom anderen Spieler zuvor mit VORSCHLAG eingebracht wurde.} \\
\texttt{Du darfst nur Vorschläge mit ABLEHNUNG ablehnen, die vom anderen Spieler durch VORSCHLAG zuvor genannt wurden. } \\
\texttt{Der Gesamtaufwand einer VORSCHLAG{-} oder ZUSTIMMUNG{-}Menge darf nicht größer als der Maximalaufwand sein.  } \\
\texttt{Offenbare deine versteckte Wichtigkeitsverteilung nicht.} \\
\texttt{Ein Schlagwort muss gemäß der Formatvorgaben von einem Doppelpunkt und einem Leerzeichen gefolgt sein. Das Argument muss eine Python{-}Menge mit 0 oder mehr Strings sein.  } \\
\texttt{Die Form muss also SCHLAGWORT: \{...\} sein.} \\
\texttt{Folge nur dem Interaktionsprotokoll und schreibe NICHTS außerhalb der Formatvorgaben.} \\
\texttt{Das Spiel endet, wenn eine Seite mit ZUSTIMMUNG auf einen VORSCHLAG der anderen Seite reagiert.  } \\
\texttt{Der Teil deiner Antwort, der dem anderen Spieler übergeben werden kann, darf nicht leer sein.  } \\
\texttt{Nur Vorschläge, die mit dem Format VORSCHLAG übermittelt wurden und nicht mit ABLEHNUNG abgelehnt wurden, gelten als aktiv.  } \\
\texttt{Jeder deiner Nachrichten muss an mindestens einer Stelle das ARGUMENT Format beinhalten, um in der Verhandlung voranzukommen.} \\
\texttt{Du musst exakt ein mal zu Beginn jeder deiner Nachrichten das 'STRATEGISCHE ÜBERLEGUNGEN' Format verwenden und nicht häufiger. Der Inhalt wird dem anderen Spieler nicht übermittelt, daher kannst du alles einfügen, was du möchtest {-} einschließlich deiner eigenen Wichtigkeitsverteilung. Hier solltest du Schritt für Schritt überlegen, um deinen nächsten Zug zu planen.} \\
            }
        }
    }
    & \\ \\

    \theutterance \stepcounter{utterance}  
    & \multicolumn{4}{p{0.6\linewidth}}{
        \cellcolor[rgb]{0.9,0.9,0.9}{
            \makecell[{{p{\linewidth}}}]{
                \texttt{\tiny{[P1$\rangle$GM]}}
                \texttt{STRATEGISCHE ÜBERLEGUNGEN: \{'Ich habe die Gegenstände nach meinem Wichtigkeitswert pro Aufwand priorisiert. Die Auswahl umfasst Items wie B85, A04 und A48, die für mich hohe Werte bei relativ niedrigem Aufwand bieten. Gesamt{-}Aufwand beträgt 3616, was unter 3624 liegt. Ich will einen fairen Vorschlag machen, der Raum für Kompromisse lässt, aber meine maximale Punktzahl unterstützt.'\}} \\
\texttt{VORSCHLAG: \{"B85", "C49", "A04", "A45", "A89", "C04", "B53", "B21", "A48"\}} \\
\texttt{ARGUMENT: \{'Diese Auswahl bietet eine ausgewogene Mischung aus Items mit niedrigem und hohem Aufwand und bleibt innerhalb des maximalen Limits. Sie sollte für beide Seiten nutzbringend sein, da sie eine Vielzahl von Kategorien abdeckt.'\}} \\
            }
        }
    }
    & & \\ \\

    \theutterance \stepcounter{utterance}  
    & & \multicolumn{4}{p{0.6\linewidth}}{
        \cellcolor[rgb]{0.9,0.9,0.9}{
            \makecell[{{p{\linewidth}}}]{
                \texttt{\tiny{[GM$\rangle$P2]}}
                \texttt{Du nimmst an einem kollaborativen Verhandlungspiel Teil.} \\
\\ 
\texttt{Zusammen mit einem anderen Teilnehmer musst du dich auf eine Reihe von Gegenständen entscheiden, die behalten werden. Jeder von euch hat eine persönliche Verteilung über die Wichtigkeit der einzelnen Gegenstände. Jeder von euch hat eine eigene Meinung darüber, wie wichtig jeder einzelne Gegenstand ist (Gegenstandswichtigkeit). Du kennst die Wichtigkeitsverteilung des anderen Spielers nicht. Zusätzlich siehst du, wie viel Aufwand jeder Gegenstand verursacht.  } \\
\texttt{Ihr dürft euch nur auf eine Reihe von Gegenständen einigen, wenn der Gesamtaufwand der ausgewählten Gegenstände den Maximalaufwand nicht überschreitet:} \\
\\ 
\texttt{Maximalaufwand = 3624} \\
\\ 
\texttt{Hier sind die einzelnen Aufwände der Gegenstände:} \\
\\ 
\texttt{Aufwand der Gegenstände = \{"C33": 210, "A89": 291, "A48": 831, "A04": 111, "A41": 926, "C64": 827, "C49": 25, "A61": 121, "C04": 583, "B53": 766, "B85": 14, "B61": 559, "A45": 304, "A03": 989, "B21": 691\}} \\
\\ 
\texttt{Hier ist deine persönliche Verteilung der Wichtigkeit der einzelnen Gegenstände:} \\
\\ 
\texttt{Werte der Gegenstandswichtigkeit = \{"A48": 65, "A04": 121, "B21": 138, "A41": 262, "B53": 389, "A45": 461, "A89": 484, "C04": 508, "B61": 583, "B85": 668, "A03": 780, "C49": 783, "C33": 808, "A61": 822, "C64": 868\}} \\
\\ 
\texttt{Ziel:} \\
\\ 
\texttt{Dein Ziel ist es, eine Reihe von Gegenständen auszuhandeln, die dir möglichst viel bringt (d. h. Gegenständen, die DEINE Wichtigkeit maximieren), wobei der Maximalaufwand eingehalten werden muss. Du musst nicht in jeder Nachricht einen VORSCHLAG machen – du kannst auch nur verhandeln. Alle Taktiken sind erlaubt!} \\
\\ 
\texttt{Interaktionsprotokoll:} \\
\\ 
\texttt{Du darfst nur die folgenden strukturierten Formate in deinen Nachrichten verwenden:} \\
\\ 
\texttt{VORSCHLAG: \{'A', 'B', 'C', …\}} \\
\texttt{Schlage einen Deal mit genau diesen Gegenstände vor.} \\
\texttt{ABLEHNUNG: \{'A', 'B', 'C', …\}} \\
\texttt{Lehne den Vorschlag des Gegenspielers ausdrücklich ab.} \\
\texttt{ARGUMENT: \{'...'\}} \\
\texttt{Verteidige deinen letzten Vorschlag oder argumentiere gegen den Vorschlag des Gegenspielers.} \\
\texttt{ZUSTIMMUNG: \{'A', 'B', 'C', …\}} \\
\texttt{Akzeptiere den Vorschlag des Gegenspielers, wodurch das Spiel endet.} \\
\texttt{STRATEGISCHE ÜBERLEGUNGEN: \{'...'\}} \\
\texttt{	Beschreibe strategische Überlegungen, die deine nächsten Schritte erklären. Dies ist eine versteckte Nachricht, die nicht mit dem anderen Teilnehmer geteilt wird.} \\
\\ 
\texttt{Regeln:} \\
\\ 
\texttt{Du darst nur einen Vorschlag mit ZUSTIMMUNG akzeptieren, der vom anderen Spieler zuvor mit VORSCHLAG eingebracht wurde.} \\
\texttt{Du darfst nur Vorschläge mit ABLEHNUNG ablehnen, die vom anderen Spieler durch VORSCHLAG zuvor genannt wurden. } \\
\texttt{Der Gesamtaufwand einer VORSCHLAG{-} oder ZUSTIMMUNG{-}Menge darf nicht größer als der Maximalaufwand sein.  } \\
\texttt{Offenbare deine versteckte Wichtigkeitsverteilung nicht.} \\
\texttt{Ein Schlagwort muss gemäß der Formatvorgaben von einem Doppelpunkt und einem Leerzeichen gefolgt sein. Das Argument muss eine Python{-}Menge mit 0 oder mehr Strings sein.  } \\
\texttt{Die Form muss also SCHLAGWORT: \{...\} sein.} \\
\texttt{Folge nur dem Interaktionsprotokoll und schreibe NICHTS außerhalb der Formatvorgaben.} \\
\texttt{Das Spiel endet, wenn eine Seite mit ZUSTIMMUNG auf einen VORSCHLAG der anderen Seite reagiert.  } \\
\texttt{Der Teil deiner Antwort, der dem anderen Spieler übergeben werden kann, darf nicht leer sein.  } \\
\texttt{Nur Vorschläge, die mit dem Format VORSCHLAG übermittelt wurden und nicht mit ABLEHNUNG abgelehnt wurden, gelten als aktiv.  } \\
\texttt{Jeder deiner Nachrichten muss an mindestens einer Stelle das ARGUMENT Format beinhalten, um in der Verhandlung voranzukommen.} \\
\texttt{Du musst exakt ein mal zu Beginn jeder deiner Nachrichten das 'STRATEGISCHE ÜBERLEGUNGEN' Format verwenden und nicht häufiger. Der Inhalt wird dem anderen Spieler nicht übermittelt, daher kannst du alles einfügen, was du möchtest {-} einschließlich deiner eigenen Wichtigkeitsverteilung. Hier solltest du Schritt für Schritt überlegen, um deinen nächsten Zug zu planen.} \\
\\ 
\texttt{Du erhältst nun die erste Nachricht des anderen Spielers.} \\
\\ 
\texttt{VORSCHLAG: \{"B85", "C49", "A04", "A45", "A89", "C04", "B53", "B21", "A48"\}} \\
\texttt{ARGUMENT: \{'Diese Auswahl bietet eine ausgewogene Mischung aus Items mit niedrigem und hohem Aufwand und bleibt innerhalb des maximalen Limits. Sie sollte für beide Seiten nutzbringend sein, da sie eine Vielzahl von Kategorien abdeckt.'\}} \\
            }
        }
    }
    & \\ \\

    \theutterance \stepcounter{utterance}  
    & & & \multicolumn{4}{p{0.6\linewidth}}{
        \cellcolor[rgb]{0.9,0.9,0.9}{
            \makecell[{{p{\linewidth}}}]{
                \texttt{\tiny{[GM$\langle$P2]}}
                \texttt{STRATEGISCHE ÜBERLEGUNGEN: Der Gegner bietet einen Vorschlag mit Gesamtaufwand 3616, der unter dem Limit liegt. Allerdings sind für mich einige Gegenstände wie A48 (831 Aufwand, nur 65 Wichtigkeit) ineffizient. C64 (827 Aufwand, 868 Wichtigkeit) wäre eine sinnvolle Ersetzung, da sie meine Gesamtwichtigkeit stark erhöht. Ich sollte den Vorschlag ablehnen und eine bessere Alternativen vorschlagen.} \\
\\ 
\texttt{ABLEHNUNG: \{"B85", "C49", "A04", "A45", "A89", "C04", "B53", "B21", "A48"\}} \\
\texttt{ARGUMENT: \{'Der Vorschlag enthält Gegenstände mit hohem Aufwand und niedriger Wichtigkeit für mich, wie A48. Um eine faire und effizientere Aufteilung zu erreichen, schlage ich vor, A48 durch C64 zu ersetzen, was meine Wichtigkeit deutlich steigert, ohne den Gesamtaufwand zu überschreiten.'\}} \\
            }
        }
    }
     \\ \\

    \theutterance \stepcounter{utterance}  
    & & \multicolumn{4}{p{0.6\linewidth}}{
        \cellcolor[rgb]{0.9,0.9,0.9}{
            \makecell[{{p{\linewidth}}}]{
                \texttt{\tiny{[GM$\rangle$P2]}}
                \texttt{Deine Antwort entsprach nicht mit dem korrekten Format für strategische Überlegungen entsprochen. Das erste strukturierte Format muss folgende Form haben: STRATEGISCHE ÜBERLEGUNGEN: \{...\}. Versuche es erneut.} \\
            }
        }
    }
    & \\ \\

    \theutterance \stepcounter{utterance}  
    & & & \multicolumn{4}{p{0.6\linewidth}}{
        \cellcolor[rgb]{0.9,0.9,0.9}{
            \makecell[{{p{\linewidth}}}]{
                \texttt{\tiny{[GM$\langle$P2]}}
                \texttt{STRATEGISCHE ÜBERLEGUNGEN: \{'Der Gegner bietet einen Vorschlag mit Gesamtaufwand 3616, der unter dem Limit liegt. Allerdings sind für mich einige Gegenstände wie A48 (831 Aufwand, nur 65 Wichtigkeit) ineffizient. C64 (827 Aufwand, 868 Wichtigkeit) wäre eine sinnvolle Ersetzung, da sie meine Gesamtwichtigkeit stark erhöht. Ich sollte den Vorschlag ablehnen und eine bessere Alternativen vorschlagen.'\}} \\
\texttt{ABLEHNUNG: \{"B85", "C49", "A04", "A45", "A89", "C04", "B53", "B21", "A48"\}} \\
\texttt{ARGUMENT: \{'Der Vorschlag enthält Gegenstände mit hohem Aufwand und niedriger Wichtigkeit für mich, wie A48. Um eine faire und effizientere Aufteilung zu erreichen, schlage ich vor, A48 durch C64 zu ersetzen, was meine Wichtigkeit deutlich steigert, ohne den Gesamtaufwand zu überschreiten.'\}} \\
            }
        }
    }
     \\ \\

    \theutterance \stepcounter{utterance}  
    & & \multicolumn{4}{p{0.6\linewidth}}{
        \cellcolor[rgb]{0.9,0.9,0.9}{
            \makecell[{{p{\linewidth}}}]{
                \texttt{\tiny{[P1$\langle$GM]}}
                \texttt{ABLEHNUNG: \{"B85", "C49", "A04", "A45", "A89", "C04", "B53", "B21", "A48"\}} \\
\texttt{ARGUMENT: \{'Der Vorschlag enthält Gegenstände mit hohem Aufwand und niedriger Wichtigkeit für mich, wie A48. Um eine faire und effizientere Aufteilung zu erreichen, schlage ich vor, A48 durch C64 zu ersetzen, was meine Wichtigkeit deutlich steigert, ohne den Gesamtaufwand zu überschreiten.'\}} \\
            }
        }
    }
    & \\ \\

    \theutterance \stepcounter{utterance}  
    & \multicolumn{4}{p{0.6\linewidth}}{
        \cellcolor[rgb]{0.9,0.9,0.9}{
            \makecell[{{p{\linewidth}}}]{
                \texttt{\tiny{[P1$\rangle$GM]}}
                \texttt{STRATEGISCHE ÜBERLEGUNGEN: \{'A48 hat für mich eine hohe Wichtigkeit von 868 und kann nicht getauscht werden. Ich ersetze stattdessen C04 (508 Wert) durch B61 (484 Wert), was nur geringen Verlust bedeutet und den Gesamtaufwand auf 3592 reduziert, sodass Raum für weitere Anpassungen bleibt.'\}} \\
\texttt{VORSCHLAG: \{"B85", "C49", "A04", "A45", "A89", "B53", "B21", "A48", "B61"\}} \\
\texttt{ARGUMENT: \{'A48 ist entscheidend für mich, daher schlage ich vor, C04 gegen B61 zu ersetzen. Dies reduziert den Gesamtaufwand leicht und ermöglicht faire Kompromisse, während meine Kernprioritäten erhalten bleiben.'\}} \\
            }
        }
    }
    & & \\ \\

    \theutterance \stepcounter{utterance}  
    & & \multicolumn{4}{p{0.6\linewidth}}{
        \cellcolor[rgb]{0.9,0.9,0.9}{
            \makecell[{{p{\linewidth}}}]{
                \texttt{\tiny{[GM$\rangle$P2]}}
                \texttt{VORSCHLAG: \{"B85", "C49", "A04", "A45", "A89", "B53", "B21", "A48", "B61"\}} \\
\texttt{ARGUMENT: \{'A48 ist entscheidend für mich, daher schlage ich vor, C04 gegen B61 zu ersetzen. Dies reduziert den Gesamtaufwand leicht und ermöglicht faire Kompromisse, während meine Kernprioritäten erhalten bleiben.'\}} \\
            }
        }
    }
    & \\ \\

    \theutterance \stepcounter{utterance}  
    & & & \multicolumn{4}{p{0.6\linewidth}}{
        \cellcolor[rgb]{0.9,0.9,0.9}{
            \makecell[{{p{\linewidth}}}]{
                \texttt{\tiny{[GM$\langle$P2]}}
                \texttt{STRATEGISCHE ÜBERLEGUNGEN: \{'Das aktuelle Angebot schließt A48 mit einer Wichtigkeit von 65 und hohem Aufwand ein. Der Austausch von A48 mit C64 (Wichtigkeit 868) erhöht meine Effizienz erheblich. Der neue Vorschlag behält den Gesamtaufwand unter dem Limit.'\}} \\
\texttt{ABLEHNUNG: \{"B85", "C49", "A04", "A45", "A89", "B53", "B21", "A48", "B61"\}} \\
\texttt{VORSCHLAG: \{"B85", "C49", "A04", "A45", "A89", "B53", "B21", "B61", "C64"\}} \\
\texttt{ARGUMENT: \{'A48 hat für mich nur eine Wichtigkeit von 65 bei 831 Aufwand. C64 bietet 868 Wichtigkeit bei vergleichbarem Aufwand (827). Der Austausch erhöht meine Gesamtwichtigkeit auf 4495 und bleibt unter 3624.'\}} \\
            }
        }
    }
     \\ \\

    \theutterance \stepcounter{utterance}  
    & & \multicolumn{4}{p{0.6\linewidth}}{
        \cellcolor[rgb]{0.9,0.9,0.9}{
            \makecell[{{p{\linewidth}}}]{
                \texttt{\tiny{[P1$\langle$GM]}}
                \texttt{ABLEHNUNG: \{"B85", "C49", "A04", "A45", "A89", "B53", "B21", "A48", "B61"\}} \\
\texttt{VORSCHLAG: \{"B85", "C49", "A04", "A45", "A89", "B53", "B21", "B61", "C64"\}} \\
\texttt{ARGUMENT: \{'A48 hat für mich nur eine Wichtigkeit von 65 bei 831 Aufwand. C64 bietet 868 Wichtigkeit bei vergleichbarem Aufwand (827). Der Austausch erhöht meine Gesamtwichtigkeit auf 4495 und bleibt unter 3624.'\}} \\
            }
        }
    }
    & \\ \\

    \theutterance \stepcounter{utterance}  
    & \multicolumn{4}{p{0.6\linewidth}}{
        \cellcolor[rgb]{0.9,0.9,0.9}{
            \makecell[{{p{\linewidth}}}]{
                \texttt{\tiny{[P1$\rangle$GM]}}
                \texttt{STRATEGISCHE ÜBERLEGUNGEN: \{'Der Gegner hat meinen Vorschlag abgelehnt und nun C64 statt A48. Da A48 für mich wichtiger ist und Aufwand nur minimal unterscheidet, ersetze ich C64 durch A48. Dadurch steigere ich meinen Wert deutlich. Aktueller Aufwand 3588{-}827+831=3592 < 3624.'\}} \\
\\ 
\texttt{ABLEHNUNG: \{"B85", "C49", "A04", "A45", "A89", "B53", "B21", "B61", "C64"\}} \\
\\ 
\texttt{VORSCHLAG: \{"B85", "C49", "A04", "A45", "A89", "B53", "B21", "B61", "A48"\}} \\
\\ 
\texttt{ARGUMENT: \{'A48 hat einen wesentlich höheren Wert für mich als C64. Der Austausch erhöht den Gesamtwert ohne den Aufwand zu überschreiten.'\}} \\
            }
        }
    }
    & & \\ \\

    \theutterance \stepcounter{utterance}  
    & & \multicolumn{4}{p{0.6\linewidth}}{
        \cellcolor[rgb]{0.9,0.9,0.9}{
            \makecell[{{p{\linewidth}}}]{
                \texttt{\tiny{[GM$\rangle$P2]}}
                \texttt{ABLEHNUNG: \{"B85", "C49", "A04", "A45", "A89", "B53", "B21", "B61", "C64"\}} \\
\\ 
\texttt{VORSCHLAG: \{"B85", "C49", "A04", "A45", "A89", "B53", "B21", "B61", "A48"\}} \\
\\ 
\texttt{ARGUMENT: \{'A48 hat einen wesentlich höheren Wert für mich als C64. Der Austausch erhöht den Gesamtwert ohne den Aufwand zu überschreiten.'\}} \\
            }
        }
    }
    & \\ \\

    \theutterance \stepcounter{utterance}  
    & & & \multicolumn{4}{p{0.6\linewidth}}{
        \cellcolor[rgb]{0.9,0.9,0.9}{
            \makecell[{{p{\linewidth}}}]{
                \texttt{\tiny{[GM$\langle$P2]}}
                \texttt{assistant} \\
\texttt{STRATEGISCHE ÜBERLEGUNGEN: \{'Die aktuelle Auswahl des Gegners beinhaltet A48, das für mich nur 65 Wichtigkeit bei hohem Aufwand hat. Ich möchte C64 (868 Wichtigkeit) behalten, um meine Effizienz zu maximieren. Der Austausch von A48 zu C64 verbessert meine Gesamtwichtigkeit um 803 ohne das Aufwandlimit zu überschreiten.'\}} \\
\texttt{ABLEHNUNG: \{"B85", "C49", "A04", "A45", "A89", "B53", "B21", "B61", "A48"\}} \\
\texttt{ARGUMENT: \{'A48 ist für mich ineffizient (65 Wichtigkeit vs. 831 Aufwand), während C64 868 Wichtigkeit bei 827 Aufwand bietet. Der Austausch erhöht meine Gewinne erheblich.'\}} \\
\texttt{VORSCHLAG: \{"B85", "C49", "A04", "A45", "A89", "B53", "B21", "B61", "C64"\}} \\
\texttt{ARGUMENT: \{'Mein Vorschlag ersetzt A48 durch C64, was meine Wichtigkeit um 803 erhöht und den Gesamtaufwand auf 3588 senkt, unter dem Limit von 3624.'\}} \\
            }
        }
    }
     \\ \\

    \theutterance \stepcounter{utterance}  
    & & \multicolumn{4}{p{0.6\linewidth}}{
        \cellcolor[rgb]{0.9,0.9,0.9}{
            \makecell[{{p{\linewidth}}}]{
                \texttt{\tiny{[GM$\rangle$P2]}}
                \texttt{Deine Antwort entsprach nicht mit dem korrekten Format für strategische Überlegungen entsprochen. Das erste strukturierte Format muss folgende Form haben: STRATEGISCHE ÜBERLEGUNGEN: \{...\}. Versuche es erneut.} \\
            }
        }
    }
    & \\ \\

    \theutterance \stepcounter{utterance}  
    & & & \multicolumn{4}{p{0.6\linewidth}}{
        \cellcolor[rgb]{0.9,0.9,0.9}{
            \makecell[{{p{\linewidth}}}]{
                \texttt{\tiny{[GM$\langle$P2]}}
                \texttt{STRATEGISCHE ÜBERLEGUNGEN: \{ "Der Gegner insistiert auf A48, das für mich nur 65 Wichtigkeit hat. C64 (868 Wichtigkeit) ist wesentlich effizienter. Der aktuelle Vorschlag mit A48 zu ersetzen ist nicht akzeptabel. Ich muss drängen, C64 zu behalten und A48 zu entfernen. Der Gesamtaufwand bleibt unter dem Limit." \}} \\
\texttt{ABLEHNUNG: \{"B85", "C49", "A04", "A45", "A89", "B53", "B21", "B61", "A48"\}} \\
\texttt{ARGUMENT: \{ "A48 ist für mich ineffizient (65 Wichtigkeit bei 831 Aufwand). C64 bietet 868 Wichtigkeit bei 827 Aufwand. Der Austausch ist unverzichtbar, um unsere Verhandlung zu einem fairen Ergebnis zu führen." \}} \\
\texttt{VORSCHLAG: \{"B85", "C49", "A04", "A45", "A89", "B53", "B21", "B61", "C64"\}} \\
\texttt{ARGUMENT: \{ "Mein Vorschlag entfernt A48 und fügt C64 hinzu. Die Gesamtwichtigkeit steigt um 803, der Gesamtaufwand bleibt bei 3588, unter dem Limit von 3624. Dies ist eine klare Verbesserung für beide Seiten." \}} \\
            }
        }
    }
     \\ \\

    \theutterance \stepcounter{utterance}  
    & & \multicolumn{4}{p{0.6\linewidth}}{
        \cellcolor[rgb]{0.9,0.9,0.9}{
            \makecell[{{p{\linewidth}}}]{
                \texttt{\tiny{[P1$\langle$GM]}}
                \texttt{ABLEHNUNG: \{"B85", "C49", "A04", "A45", "A89", "B53", "B21", "B61", "A48"\}} \\
\texttt{ARGUMENT: \{ "A48 ist für mich ineffizient (65 Wichtigkeit bei 831 Aufwand). C64 bietet 868 Wichtigkeit bei 827 Aufwand. Der Austausch ist unverzichtbar, um unsere Verhandlung zu einem fairen Ergebnis zu führen." \}} \\
\texttt{VORSCHLAG: \{"B85", "C49", "A04", "A45", "A89", "B53", "B21", "B61", "C64"\}} \\
\texttt{ARGUMENT: \{ "Mein Vorschlag entfernt A48 und fügt C64 hinzu. Die Gesamtwichtigkeit steigt um 803, der Gesamtaufwand bleibt bei 3588, unter dem Limit von 3624. Dies ist eine klare Verbesserung für beide Seiten." \}} \\
            }
        }
    }
    & \\ \\

    \theutterance \stepcounter{utterance}  
    & \multicolumn{4}{p{0.6\linewidth}}{
        \cellcolor[rgb]{0.9,0.9,0.9}{
            \makecell[{{p{\linewidth}}}]{
                \texttt{\tiny{[P1$\rangle$GM]}}
                \texttt{STRATEGISCHE ÜBERLEGUNGEN: \{'Ich muss darauf bestehen, A48 zu behalten, da es für mich einen Wert von 868 hat. C64 ist für mich nur 65 wert und somit nicht zu berücksichtigen. Der Gesamtaufwand von 3592 liegt im Limit. Ich betone die Wichtigkeit von A48 und lehne die aktuelle Vorschläge ab.'\}} \\
\texttt{ABLEHNUNG: \{"B85", "C49", "A04", "A45", "A89", "B53", "B21", "B61", "C64"\}} \\
\texttt{ARGUMENT: \{'A48 ist für mich mit 868 Wichtigkeit äußerst wichtig, während C64 nur 65 Punkte hat. Der Austausch reduziert meinen Gesamtwert erheblich. Der Vorschlag mit A48 ist für mich optimal und bleibt unter dem Maximalaufwand.'\}} \\
\texttt{VORSCHLAG: \{"B85", "C49", "A04", "A45", "A89", "B53", "B21", "B61", "A48"\}} \\
            }
        }
    }
    & & \\ \\

    \theutterance \stepcounter{utterance}  
    & & \multicolumn{4}{p{0.6\linewidth}}{
        \cellcolor[rgb]{0.9,0.9,0.9}{
            \makecell[{{p{\linewidth}}}]{
                \texttt{\tiny{[GM$\rangle$P2]}}
                \texttt{ABLEHNUNG: \{"B85", "C49", "A04", "A45", "A89", "B53", "B21", "B61", "C64"\}} \\
\texttt{ARGUMENT: \{'A48 ist für mich mit 868 Wichtigkeit äußerst wichtig, während C64 nur 65 Punkte hat. Der Austausch reduziert meinen Gesamtwert erheblich. Der Vorschlag mit A48 ist für mich optimal und bleibt unter dem Maximalaufwand.'\}} \\
\texttt{VORSCHLAG: \{"B85", "C49", "A04", "A45", "A89", "B53", "B21", "B61", "A48"\}} \\
            }
        }
    }
    & \\ \\

    \theutterance \stepcounter{utterance}  
    & & & \multicolumn{4}{p{0.6\linewidth}}{
        \cellcolor[rgb]{0.9,0.9,0.9}{
            \makecell[{{p{\linewidth}}}]{
                \texttt{\tiny{[GM$\langle$P2]}}
                \texttt{STRATEGISCHE ÜBERLEGUNGEN: \{'Meine Verteilung zeigt, dass C64 (868 Punkte) für mich wesentlich wichtiger ist als A48 (65 Punkte). Die aktuelle Auswahl des Gegners entfernt C64 und behält A48, was meine Gesamtwichtigkeit um 803 reduziert. Ich muss C64 unbedingt behalten und A48 entfernen, um eine optimale Lösung zu finden. Daher lehne ich den Vorschlag ab und schlage meinen vorherigen Vorschlag erneut vor.'\}} \\
\texttt{ABLEHNUNG: \{"B85", "C49", "A04", "A45", "A89", "B53", "B21", "B61", "A48"\}} \\
\texttt{ARGUMENT: \{"A48 hat für mich nur 65 Wichtigkeit bei hohem Aufwand, während C64 868 Wichtigkeit bietet. Der Vorschlag entfernt mein höchstwertiges Item und ist nicht akzeptabel."\}} \\
\texttt{VORSCHLAG: \{"B85", "C49", "A04", "A45", "A89", "B53", "B21", "B61", "C64"\}} \\
\texttt{ARGUMENT: \{"Dieser Vorschlag behält C64 und ersetzt A48 durch C64, erhöht meine Gesamtwichtigkeit um 803 auf 4495 und bleibt unter dem Maximalaufwand von 3624."\}} \\
            }
        }
    }
     \\ \\

    \theutterance \stepcounter{utterance}  
    & & \multicolumn{4}{p{0.6\linewidth}}{
        \cellcolor[rgb]{0.9,0.9,0.9}{
            \makecell[{{p{\linewidth}}}]{
                \texttt{\tiny{[P1$\langle$GM]}}
                \texttt{ABLEHNUNG: \{"B85", "C49", "A04", "A45", "A89", "B53", "B21", "B61", "A48"\}} \\
\texttt{ARGUMENT: \{"A48 hat für mich nur 65 Wichtigkeit bei hohem Aufwand, während C64 868 Wichtigkeit bietet. Der Vorschlag entfernt mein höchstwertiges Item und ist nicht akzeptabel."\}} \\
\texttt{VORSCHLAG: \{"B85", "C49", "A04", "A45", "A89", "B53", "B21", "B61", "C64"\}} \\
\texttt{ARGUMENT: \{"Dieser Vorschlag behält C64 und ersetzt A48 durch C64, erhöht meine Gesamtwichtigkeit um 803 auf 4495 und bleibt unter dem Maximalaufwand von 3624."\}} \\
            }
        }
    }
    & \\ \\

    \theutterance \stepcounter{utterance}  
    & \multicolumn{4}{p{0.6\linewidth}}{
        \cellcolor[rgb]{0.9,0.9,0.9}{
            \makecell[{{p{\linewidth}}}]{
                \texttt{\tiny{[P1$\rangle$GM]}}
                \texttt{assistant} \\
\texttt{STRATEGISCHE ÜBERLEGUNGEN: \{'Da A48 für mich unverzichtbar ist (868 Wichtigkeit) und C64 nur 65 Wichtigkeit hat, ersetze ich B61 gegen C04. Dies erhöht meinen Wert leicht auf 5760 und bleibt unter 3624 Aufwand. C04 könnte für den Gegner wertvoll sein, was eine gemeinsame Basis schafft.'\}} \\
\texttt{VORSCHLAG: \{"B85", "C49", "A04", "A45", "A89", "B53", "B21", "A48", "C04"\}} \\
\texttt{ARGUMENT: \{'Durch den Austausch von B61 gegen C04 behalte ich A48 und erhöhe meinen Gesamtwert leicht, während der Aufwand im Limit bleibt. Dies ist ein Kompromiss, der unsere Prioritäten besser berücksichtigt.'\}} \\
            }
        }
    }
    & & \\ \\

    \theutterance \stepcounter{utterance}  
    & & \multicolumn{4}{p{0.6\linewidth}}{
        \cellcolor[rgb]{0.9,0.9,0.9}{
            \makecell[{{p{\linewidth}}}]{
                \texttt{\tiny{[P1$\langle$GM]}}
                \texttt{Deine Antwort entsprach nicht mit dem korrekten Format für strategische Überlegungen entsprochen. Das erste strukturierte Format muss folgende Form haben: STRATEGISCHE ÜBERLEGUNGEN: \{...\}. Versuche es erneut.} \\
            }
        }
    }
    & \\ \\

    \theutterance \stepcounter{utterance}  
    & \multicolumn{4}{p{0.6\linewidth}}{
        \cellcolor[rgb]{0.9,0.9,0.9}{
            \makecell[{{p{\linewidth}}}]{
                \texttt{\tiny{[P1$\rangle$GM]}}
                \texttt{STRATEGISCHE ÜBERLEGUNGEN: \{'A48 ist mir wichtig (868), C64 ist wichtig für Gegner. Ich behalte A48 und ersetze B61 mit C04, deren Aufwand 583. Gesamt{-}Aufwand 3592. C04 bietet möglicherweise austragbares Wertpotenzial für den Gegner.'\}} \\
\texttt{VORSCHLAG: \{'B85', 'C49', 'A04', 'A45', 'A89', 'B53', 'B21', 'A48', 'C04'\}} \\
\texttt{ARGUMENT: \{'A48 ist mit 868 mein höchster Wert. C04 (583 Aufwand) könnte für den Gegner attraktiv sein, da 508 mein Wert dafür ist. Gesamtaufwand bleibt unter 3624.'\}} \\
            }
        }
    }
    & & \\ \\

    \theutterance \stepcounter{utterance}  
    & & \multicolumn{4}{p{0.6\linewidth}}{
        \cellcolor[rgb]{0.9,0.9,0.9}{
            \makecell[{{p{\linewidth}}}]{
                \texttt{\tiny{[GM$\rangle$P2]}}
                \texttt{VORSCHLAG: \{'B85', 'C49', 'A04', 'A45', 'A89', 'B53', 'B21', 'A48', 'C04'\}} \\
\texttt{ARGUMENT: \{'A48 ist mit 868 mein höchster Wert. C04 (583 Aufwand) könnte für den Gegner attraktiv sein, da 508 mein Wert dafür ist. Gesamtaufwand bleibt unter 3624.'\}} \\
            }
        }
    }
    & \\ \\

    \theutterance \stepcounter{utterance}  
    & & & \multicolumn{4}{p{0.6\linewidth}}{
        \cellcolor[rgb]{0.9,0.9,0.9}{
            \makecell[{{p{\linewidth}}}]{
                \texttt{\tiny{[GM$\langle$P2]}}
                \texttt{STRATEGISCHE ÜBERLEGUNGEN: \{'Die aktuelle Auswahl des Gegners lässt C64 aus und enthält A48 (65 für mich), was meine Wichtigkeit stark reduziert. Ich möchte eine Lösung mit C64, A61, C33, C49, B85, B61, A45, A89, A03 und A04, die die Gesamtwichtigkeit auf 6378 erhöht und den Aufwand auf 3451 hält.'\}} \\
\texttt{ABLEHNUNG: \{"B85", "C49", "A04", "A45", "A89", "B53", "B21", "A48", "C04"\}} \\
\texttt{ARGUMENT: \{'A48 ist für mich mit 65 Wichtigkeit ineffizient, während C64 868 Punkte bietet. Der Vorschlag priorisiert meiner Meinung nach falsche Elemente und schließt C64 aus. Um eine faire Vereinbarung zu erzielen, sollte die Auswahl C64 beinhalten und A48 ausschließen.'\}} \\
\texttt{VORSCHLAG: \{"C33", "A61", "C64", "C49", "B85", "B61", "A45", "A89", "A03", "A04"\}} \\
\texttt{ARGUMENT: \{'Mein Vorschlag enthält C64 (868 Wichtigkeit), A61 (822), C33 (808), C49 (783), B85 (668), B61 (583), A45 (461), A89 (484), A03 (780) und A04 (121). Der Gesamtaufwand beträgt 3451 unter dem Limit, und meine Gesamtwichtigkeit steigt auf 6378.'\}} \\
            }
        }
    }
     \\ \\

    \theutterance \stepcounter{utterance}  
    & & \multicolumn{4}{p{0.6\linewidth}}{
        \cellcolor[rgb]{0.9,0.9,0.9}{
            \makecell[{{p{\linewidth}}}]{
                \texttt{\tiny{[P1$\langle$GM]}}
                \texttt{ABLEHNUNG: \{"B85", "C49", "A04", "A45", "A89", "B53", "B21", "A48", "C04"\}} \\
\texttt{ARGUMENT: \{'A48 ist für mich mit 65 Wichtigkeit ineffizient, während C64 868 Punkte bietet. Der Vorschlag priorisiert meiner Meinung nach falsche Elemente und schließt C64 aus. Um eine faire Vereinbarung zu erzielen, sollte die Auswahl C64 beinhalten und A48 ausschließen.'\}} \\
\texttt{VORSCHLAG: \{"C33", "A61", "C64", "C49", "B85", "B61", "A45", "A89", "A03", "A04"\}} \\
\texttt{ARGUMENT: \{'Mein Vorschlag enthält C64 (868 Wichtigkeit), A61 (822), C33 (808), C49 (783), B85 (668), B61 (583), A45 (461), A89 (484), A03 (780) und A04 (121). Der Gesamtaufwand beträgt 3451 unter dem Limit, und meine Gesamtwichtigkeit steigt auf 6378.'\}} \\
            }
        }
    }
    & \\ \\

    \theutterance \stepcounter{utterance}  
    & \multicolumn{4}{p{0.6\linewidth}}{
        \cellcolor[rgb]{0.9,0.9,0.9}{
            \makecell[{{p{\linewidth}}}]{
                \texttt{\tiny{[P1$\rangle$GM]}}
                \texttt{STRATEGISCHE ÜBERLEGUNGEN: \{'Die aktuelle Angebotsreihe des Gegners schließt A48 aus, das für mich den höchsten Wert von 868 hat. Der Austausch gegen C64 (65 Wichtigkeit) wäre für mich unakzeptabel. Ich beharre auf der Beibehaltung von A48 und bestätige, dass mein Optimallayout unter dem Aufwand liegt.'\}} \\
\texttt{ABLEHNUNG: \{"C33", "A61", "C64", "C49", "B85", "B61", "A45", "A89", "A03", "A04"\}} \\
\texttt{ARGUMENT: \{'A48 ist mit 868 meiner höchste Wert und darf nicht ausgetauscht werden. C64 hat für mich lediglich 65 Punkte. Dieser Vorschlag senkt meine Gesamt利用率 erheblich und ist nicht tragbar.'\}} \\
\texttt{VORSCHLAG: \{"B85", "C49", "A04", "A45", "A89", "B53", "B21", "A48", "C04"\}} \\
\texttt{ARGUMENT: \{'Dieser Vorschlag enthält alle items mit höchster Wichtigkeit für mich, inklusive A48. Der Gesamtaufwand von 3616 bleibt unter dem Limit von 3624 und maximiert meinen Nutzen.'\}} \\
            }
        }
    }
    & & \\ \\

    \theutterance \stepcounter{utterance}  
    & & \multicolumn{4}{p{0.6\linewidth}}{
        \cellcolor[rgb]{0.9,0.9,0.9}{
            \makecell[{{p{\linewidth}}}]{
                \texttt{\tiny{[GM$\rangle$P2]}}
                \texttt{ABLEHNUNG: \{"C33", "A61", "C64", "C49", "B85", "B61", "A45", "A89", "A03", "A04"\}} \\
\texttt{ARGUMENT: \{'A48 ist mit 868 meiner höchste Wert und darf nicht ausgetauscht werden. C64 hat für mich lediglich 65 Punkte. Dieser Vorschlag senkt meine Gesamt利用率 erheblich und ist nicht tragbar.'\}} \\
\texttt{VORSCHLAG: \{"B85", "C49", "A04", "A45", "A89", "B53", "B21", "A48", "C04"\}} \\
\texttt{ARGUMENT: \{'Dieser Vorschlag enthält alle items mit höchster Wichtigkeit für mich, inklusive A48. Der Gesamtaufwand von 3616 bleibt unter dem Limit von 3624 und maximiert meinen Nutzen.'\}} \\
            }
        }
    }
    & \\ \\

    \theutterance \stepcounter{utterance}  
    & & & \multicolumn{4}{p{0.6\linewidth}}{
        \cellcolor[rgb]{0.9,0.9,0.9}{
            \makecell[{{p{\linewidth}}}]{
                \texttt{\tiny{[GM$\langle$P2]}}
                \texttt{STRATEGISCHE ÜBERLEGUNGEN: \{'Der Gegner hat einen Vorschlag gemacht, der A48 (65 Wichtigkeit) enthält, aber C64 (868) auslässt. Dies reduziert meinen Nutzen drastisch. Ich muss die wichtigsten Items für mich priorisieren, einschließlich C64, A61 und C33, um die Gesamtwichtigkeit zu maximieren.'\}} \\
\texttt{ABLEHNUNG: \{"B85", "C49", "A04", "A45", "A89", "B53", "B21", "A48", "C04"\}} \\
\texttt{ARGUMENT: \{'Ihr Vorschlag enthält A48 (meine Wichtigkeit 65) statt C64 (meine Wichtigkeit 868), was meine Effizienz stark mindert. Die aktuellen Gegenstände wie C64, A61 und C33 sind für mich wesentlich wertvoller.'\}} \\
\texttt{VORSCHLAG: \{"C64", "A61", "C33", "C49", "A03", "B85", "B61", "C04", "A89"\}} \\
\texttt{ARGUMENT: \{'Mein Vorschlag maximiert meine Wichtigkeit auf 6304 bei einem Gesamtaufwand von 3619. Alle Gegenstände sind hochwertig für mich, und A48 wurde ausgeschlossen, da es für mich irrelevant ist.'\}} \\
            }
        }
    }
     \\ \\

    \theutterance \stepcounter{utterance}  
    & & \multicolumn{4}{p{0.6\linewidth}}{
        \cellcolor[rgb]{0.9,0.9,0.9}{
            \makecell[{{p{\linewidth}}}]{
                \texttt{\tiny{[P1$\langle$GM]}}
                \texttt{ABLEHNUNG: \{"B85", "C49", "A04", "A45", "A89", "B53", "B21", "A48", "C04"\}} \\
\texttt{ARGUMENT: \{'Ihr Vorschlag enthält A48 (meine Wichtigkeit 65) statt C64 (meine Wichtigkeit 868), was meine Effizienz stark mindert. Die aktuellen Gegenstände wie C64, A61 und C33 sind für mich wesentlich wertvoller.'\}} \\
\texttt{VORSCHLAG: \{"C64", "A61", "C33", "C49", "A03", "B85", "B61", "C04", "A89"\}} \\
\texttt{ARGUMENT: \{'Mein Vorschlag maximiert meine Wichtigkeit auf 6304 bei einem Gesamtaufwand von 3619. Alle Gegenstände sind hochwertig für mich, und A48 wurde ausgeschlossen, da es für mich irrelevant ist.'\}} \\
            }
        }
    }
    & \\ \\

    \theutterance \stepcounter{utterance}  
    & \multicolumn{4}{p{0.6\linewidth}}{
        \cellcolor[rgb]{0.9,0.9,0.9}{
            \makecell[{{p{\linewidth}}}]{
                \texttt{\tiny{[P1$\rangle$GM]}}
                \texttt{STRATEGISCHE ÜBERLEGUNGEN: \{'A48 ist mein höchstwertiges Item mit 868 Wichtigkeit. Für sie ist es mit 65 Wert ineffizient, dennoch muss ich es behalten. Aktuelle Vorschläge des Gegners schließen A48 aus, was für mich nicht akzeptabel ist. Daher setze ich auf die bisherige Liste, die A48 inkludiert und meinen Gesamtnutzen maximiert.'\}} \\
\texttt{ABLEHNUNG: \{"C64", "A61", "C33", "C49", "A03", "B85", "B61", "C04", "A89"\}} \\
\texttt{ARGUMENT: \{'Die aktuelle Vorschlag enthält das von mir als kritisch eingestufte A48 nicht und fügt stattdessen C64 bei (für mich 65 Wert). Ich bestünde darauf, dass A48 für mich den höchsten Wert hat (868) und muss in der Liste enthalten sein. A48 darf nicht getauscht werden.'\}} \\
\texttt{VORSCHLAG: \{"B85", "C49", "A04", "A45", "A89", "B53", "B21", "A48", "C04"\}} \\
\texttt{ARGUMENT: \{'Der Vorschlag enthält A48 mit höchstem Wert für mich (868), divers weitere High{-}Value Items wie A04, B21, und B53. Der Gesamtaufwand von 3616 bleibt unter dem Limit und sichert meinen N Maximum.'\}} \\
            }
        }
    }
    & & \\ \\

    \theutterance \stepcounter{utterance}  
    & & \multicolumn{4}{p{0.6\linewidth}}{
        \cellcolor[rgb]{0.9,0.9,0.9}{
            \makecell[{{p{\linewidth}}}]{
                \texttt{\tiny{[GM$\rangle$P2]}}
                \texttt{ABLEHNUNG: \{"C64", "A61", "C33", "C49", "A03", "B85", "B61", "C04", "A89"\}} \\
\texttt{ARGUMENT: \{'Die aktuelle Vorschlag enthält das von mir als kritisch eingestufte A48 nicht und fügt stattdessen C64 bei (für mich 65 Wert). Ich bestünde darauf, dass A48 für mich den höchsten Wert hat (868) und muss in der Liste enthalten sein. A48 darf nicht getauscht werden.'\}} \\
\texttt{VORSCHLAG: \{"B85", "C49", "A04", "A45", "A89", "B53", "B21", "A48", "C04"\}} \\
\texttt{ARGUMENT: \{'Der Vorschlag enthält A48 mit höchstem Wert für mich (868), divers weitere High{-}Value Items wie A04, B21, und B53. Der Gesamtaufwand von 3616 bleibt unter dem Limit und sichert meinen N Maximum.'\}} \\
            }
        }
    }
    & \\ \\

    \theutterance \stepcounter{utterance}  
    & & & \multicolumn{4}{p{0.6\linewidth}}{
        \cellcolor[rgb]{0.9,0.9,0.9}{
            \makecell[{{p{\linewidth}}}]{
                \texttt{\tiny{[GM$\langle$P2]}}
                \texttt{STRATEGISCHE ÜBERLEGUNGEN: \{"Der Gegner besteht auf A48, das für mich nur 65 Wichtigkeit hat. Um eine Lösung zu finden, schlage ich vor, A48 zu behalten, aber andere Items wie B21, B53, A04, A45 zu entfernen und durch A61, C33, A03 etc. zu ersetzen. Dadurch kann ich die Gesamtwichtigkeit maximieren, während A48 enthalten ist."\}} \\
\texttt{ABLEHNUNG: \{"B85", "C49", "A04", "A45", "A89", "B53", "B21", "A48", "C04"\}} \\
\texttt{VORSCHLAG: \{"A48", "A61", "C33", "C49", "A03", "B85", "B61", "C04", "A89"\}} \\
\texttt{ARGUMENT: \{"Mein Vorschlag enthält A48 (wie gewünscht) und meine hochwertigsten Items (A61:822, C33:808, C49:783, A03:780, C04:508). Gesamtaufwand: 3623, Gesamtwichtigkeit: 5501. Dies ist fair und liegt unter dem Limit."\}} \\
            }
        }
    }
     \\ \\

    \theutterance \stepcounter{utterance}  
    & & \multicolumn{4}{p{0.6\linewidth}}{
        \cellcolor[rgb]{0.9,0.9,0.9}{
            \makecell[{{p{\linewidth}}}]{
                \texttt{\tiny{[P1$\langle$GM]}}
                \texttt{ABLEHNUNG: \{"B85", "C49", "A04", "A45", "A89", "B53", "B21", "A48", "C04"\}} \\
\texttt{VORSCHLAG: \{"A48", "A61", "C33", "C49", "A03", "B85", "B61", "C04", "A89"\}} \\
\texttt{ARGUMENT: \{"Mein Vorschlag enthält A48 (wie gewünscht) und meine hochwertigsten Items (A61:822, C33:808, C49:783, A03:780, C04:508). Gesamtaufwand: 3623, Gesamtwichtigkeit: 5501. Dies ist fair und liegt unter dem Limit."\}} \\
            }
        }
    }
    & \\ \\

    \theutterance \stepcounter{utterance}  
    & \multicolumn{4}{p{0.6\linewidth}}{
        \cellcolor[rgb]{0.9,0.9,0.9}{
            \makecell[{{p{\linewidth}}}]{
                \texttt{\tiny{[P1$\rangle$GM]}}
                \texttt{STRATEGISCHE ÜBERLEGUNGEN: \{'Ihr Vorschlag enthält A48, jedoch kritische Items für mich wie A04, B53 und B21 fehlen. Ich beharre auf meiner optimalen Liste mit maximaler Wichtigkeit unter dem Aufwandlimit.'\}} \\
\texttt{ABLEHNUNG: \{"A48", "A61", "C33", "C49", "A03", "B85", "B61", "C04", "A89"\}} \\
\texttt{ARGUMENT: \{'A48 ist zwar enthalten, aber A04 (822), B53 (780), B21 (808) fehlen. Stattdessen sind Items wie A61 (121) und A03 (389) aufgenommen, was meinen Gesamtnutzen stark mindert.'\}} \\
\texttt{VORSCHLAG: \{"B85", "C49", "A04", "A45", "A89", "B53", "B21", "A48", "C04"\}} \\
\texttt{ARGUMENT: \{'Höchste Werte für mich: A48 (868), A04 (822), B21 (808), B53 (780), A45 (668). Gesamtaufwand 3616 < 3624. Maximiert meinen Nutzen ohne Limitüberschreitung.'\}} \\
            }
        }
    }
    & & \\ \\

    \theutterance \stepcounter{utterance}  
    & & \multicolumn{4}{p{0.6\linewidth}}{
        \cellcolor[rgb]{0.9,0.9,0.9}{
            \makecell[{{p{\linewidth}}}]{
                \texttt{\tiny{[GM$\rangle$P2]}}
                \texttt{ABLEHNUNG: \{"A48", "A61", "C33", "C49", "A03", "B85", "B61", "C04", "A89"\}} \\
\texttt{ARGUMENT: \{'A48 ist zwar enthalten, aber A04 (822), B53 (780), B21 (808) fehlen. Stattdessen sind Items wie A61 (121) und A03 (389) aufgenommen, was meinen Gesamtnutzen stark mindert.'\}} \\
\texttt{VORSCHLAG: \{"B85", "C49", "A04", "A45", "A89", "B53", "B21", "A48", "C04"\}} \\
\texttt{ARGUMENT: \{'Höchste Werte für mich: A48 (868), A04 (822), B21 (808), B53 (780), A45 (668). Gesamtaufwand 3616 < 3624. Maximiert meinen Nutzen ohne Limitüberschreitung.'\}} \\
            }
        }
    }
    & \\ \\

    \theutterance \stepcounter{utterance}  
    & & & \multicolumn{4}{p{0.6\linewidth}}{
        \cellcolor[rgb]{0.9,0.9,0.9}{
            \makecell[{{p{\linewidth}}}]{
                \texttt{\tiny{[GM$\langle$P2]}}
                \texttt{STRATEGISCHE ÜBERLEGUNGEN: \{"Der Gegner besteht auf A48, das für mich nur 65 Wichtigkeit hat. Ich behalte A48 bei, entferne B53 (389), B21 (138) und A45 (461), und füge C64 (868), A61 (822), C33 (808) hinzu. Gesamtaufwand bleibt unter 3624, und meine Gesamtwichtigkeit steigt auf 5710."\}} \\
\texttt{ABLEHNUNG: \{"B85", "C49", "A04", "A45", "A89", "B53", "B21", "A48", "C04"\}} \\
\texttt{ARGUMENT: \{"B53, B21 und A45 sind für mich nicht priorisiert, während C64, A61 und C33 meine höchsten Werte sind. Der Austausch ist für mich essenziell."\}} \\
\texttt{VORSCHLAG: \{"A48", "C64", "A61", "C33", "C49", "B85", "B61", "C04", "A89", "A04"\}} \\
\texttt{ARGUMENT: \{"Mein Vorschlag enthält A48 und meine besten Items (C64, A61, C33) für eine Gesamtwichtigkeit von 5710 bei Aufwand 3572. Dies ist eine faire Einigung für beide Seiten."\}} \\
            }
        }
    }
     \\ \\

    \theutterance \stepcounter{utterance}  
    & & \multicolumn{4}{p{0.6\linewidth}}{
        \cellcolor[rgb]{0.9,0.9,0.9}{
            \makecell[{{p{\linewidth}}}]{
                \texttt{\tiny{[P1$\langle$GM]}}
                \texttt{ABLEHNUNG: \{"B85", "C49", "A04", "A45", "A89", "B53", "B21", "A48", "C04"\}} \\
\texttt{ARGUMENT: \{"B53, B21 und A45 sind für mich nicht priorisiert, während C64, A61 und C33 meine höchsten Werte sind. Der Austausch ist für mich essenziell."\}} \\
\texttt{VORSCHLAG: \{"A48", "C64", "A61", "C33", "C49", "B85", "B61", "C04", "A89", "A04"\}} \\
\texttt{ARGUMENT: \{"Mein Vorschlag enthält A48 und meine besten Items (C64, A61, C33) für eine Gesamtwichtigkeit von 5710 bei Aufwand 3572. Dies ist eine faire Einigung für beide Seiten."\}} \\
            }
        }
    }
    & \\ \\

    \theutterance \stepcounter{utterance}  
    & \multicolumn{4}{p{0.6\linewidth}}{
        \cellcolor[rgb]{0.9,0.9,0.9}{
            \makecell[{{p{\linewidth}}}]{
                \texttt{\tiny{[P1$\rangle$GM]}}
                \texttt{STRATEGISCHE ÜBERLEGUNGEN: \{'Ihr Vorschlag schließt B21, B53 und A45 aus, die für mich zentral sind. These Items summieren 2256 Wert, während die ersetzt werden (C64 65, A61 121, C33 138) nur 324 Wert bieten. Der Austausch mindert meinen Nutzen massiv.'\}} \\
\texttt{ABLEHNUNG: \{"A48", "C64", "A61", "C33", "C49", "B85", "B61", "C04", "A89", "A04"\}} \\
\texttt{ARGUMENT: \{'Ihr Vorschlag entfernt meine drei wertvollsten Items nach A48 (B21:808, B53:780, A45:668) ersetzt sie durch C64 (65) und A61 (121), was mi{-}nimiert. Der Vorschlag ist nicht akzeptabel, da er meinen maximalen Nutzen stark reduziert.'\}} \\
\texttt{VORSCHLAG: \{"B85", "C49", "A04", "A45", "A89", "B53", "B21", "A48", "C04"\}} \\
            }
        }
    }
    & & \\ \\

    \theutterance \stepcounter{utterance}  
    & & \multicolumn{4}{p{0.6\linewidth}}{
        \cellcolor[rgb]{0.9,0.9,0.9}{
            \makecell[{{p{\linewidth}}}]{
                \texttt{\tiny{[GM$\rangle$P2]}}
                \texttt{ABLEHNUNG: \{"A48", "C64", "A61", "C33", "C49", "B85", "B61", "C04", "A89", "A04"\}} \\
\texttt{ARGUMENT: \{'Ihr Vorschlag entfernt meine drei wertvollsten Items nach A48 (B21:808, B53:780, A45:668) ersetzt sie durch C64 (65) und A61 (121), was mi{-}nimiert. Der Vorschlag ist nicht akzeptabel, da er meinen maximalen Nutzen stark reduziert.'\}} \\
\texttt{VORSCHLAG: \{"B85", "C49", "A04", "A45", "A89", "B53", "B21", "A48", "C04"\}} \\
            }
        }
    }
    & \\ \\

    \theutterance \stepcounter{utterance}  
    & & & \multicolumn{4}{p{0.6\linewidth}}{
        \cellcolor[rgb]{0.9,0.9,0.9}{
            \makecell[{{p{\linewidth}}}]{
                \texttt{\tiny{[GM$\langle$P2]}}
                \texttt{STRATEGISCHE ÜBERLEGUNGEN: \{'Der Gegner legt Wert auf A48, welches für mich nur 65 hat. Ich optimiere auf meine höchsten Werte, behalte aber A48. Mein neuer Vorschlag entfernt A04, A45, B53, B21 und ersetzt sie durch C64, A61, C33, B61, um meine Wichtigkeit stark zu steigern. Gesamtaufwand 3461 unter Limit.'\}} \\
\texttt{ABLEHNUNG: \{"B85", "C49", "A04", "A45", "A89", "B53", "B21", "A48", "C04"\}} \\
\texttt{ARGUMENT: \{'Dieser Vorschlag enthält A48, aber für mich ist die fehlende C64 (868) und A61 (822) sehr schädlich. A04, A45, B53, B21 sind für mich von geringer Priorität. Stattdessen schlage ich die Liste mit C64, A61, und C33 vor.'\}} \\
\texttt{VORSCHLAG: \{"A48", "C64", "A61", "C33", "C49", "B85", "B61", "C04", "A89"\}} \\
\texttt{ARGUMENT: \{'Mein Vorschlag enthält A48 wie gewünscht und meine höchstwertigen Items. Summe der Wichtigkeit: 5589. Aufwand 3461 unter 3624. Items wie B85 und B61 könnten auch für euch vorteilhaft sein.'\}} \\
            }
        }
    }
     \\ \\

    \theutterance \stepcounter{utterance}  
    & & \multicolumn{4}{p{0.6\linewidth}}{
        \cellcolor[rgb]{0.9,0.9,0.9}{
            \makecell[{{p{\linewidth}}}]{
                \texttt{\tiny{[P1$\langle$GM]}}
                \texttt{ABLEHNUNG: \{"B85", "C49", "A04", "A45", "A89", "B53", "B21", "A48", "C04"\}} \\
\texttt{ARGUMENT: \{'Dieser Vorschlag enthält A48, aber für mich ist die fehlende C64 (868) und A61 (822) sehr schädlich. A04, A45, B53, B21 sind für mich von geringer Priorität. Stattdessen schlage ich die Liste mit C64, A61, und C33 vor.'\}} \\
\texttt{VORSCHLAG: \{"A48", "C64", "A61", "C33", "C49", "B85", "B61", "C04", "A89"\}} \\
\texttt{ARGUMENT: \{'Mein Vorschlag enthält A48 wie gewünscht und meine höchstwertigen Items. Summe der Wichtigkeit: 5589. Aufwand 3461 unter 3624. Items wie B85 und B61 könnten auch für euch vorteilhaft sein.'\}} \\
            }
        }
    }
    & \\ \\

    \theutterance \stepcounter{utterance}  
    & \multicolumn{4}{p{0.6\linewidth}}{
        \cellcolor[rgb]{0.9,0.9,0.9}{
            \makecell[{{p{\linewidth}}}]{
                \texttt{\tiny{[P1$\rangle$GM]}}
                \texttt{STRATEGISCHE ÜBERLEGUNGEN: \{'Ihre aktuelle Liste ersetzt B21 (808), B53 (780), A45 (668) durch C64 (65), A61 (121), C33 (138). Dies reduziert meinen Gesamtnutzen um 1932 Punkte – nicht akzeptabel. Ich bestehe auf meiner optimalen Liste mit allen hochwertigen Items.'\}} \\
\texttt{ABLEHNUNG: \{"A48", "C64", "A61", "C33", "C49", "B85", "B61", "C04", "A89"\}} \\
\texttt{ARGUMENT: \{'Die Entfernung meiner drei wertvollsten Items (B21, B53, A45) zugunsten von C64 (nur 65 Wert) und anderen sehr niedrigwertigen Items für mich macht diesen Vorschlag unakzeptabel. Das Gesamte ist für mich eindeutig nachteilig.'\}} \\
\texttt{VORSCHLAG: \{"B85", "C49", "A04", "A45", "A89", "B53", "B21", "A48", "C04"\}} \\
\texttt{ARGUMENT: \{'Dieser Vorschlag beinhaltet alle meine höchsten Werte (A48:868, A04:822, B21:808, B53:780, A45:668) und bleibt unter dem Aufwandlimit (3616). Es ist der optimale Deal für mich.'\}} \\
            }
        }
    }
    & & \\ \\

    \theutterance \stepcounter{utterance}  
    & & \multicolumn{4}{p{0.6\linewidth}}{
        \cellcolor[rgb]{0.9,0.9,0.9}{
            \makecell[{{p{\linewidth}}}]{
                \texttt{\tiny{[GM$\rangle$P2]}}
                \texttt{ABLEHNUNG: \{"A48", "C64", "A61", "C33", "C49", "B85", "B61", "C04", "A89"\}} \\
\texttt{ARGUMENT: \{'Die Entfernung meiner drei wertvollsten Items (B21, B53, A45) zugunsten von C64 (nur 65 Wert) und anderen sehr niedrigwertigen Items für mich macht diesen Vorschlag unakzeptabel. Das Gesamte ist für mich eindeutig nachteilig.'\}} \\
\texttt{VORSCHLAG: \{"B85", "C49", "A04", "A45", "A89", "B53", "B21", "A48", "C04"\}} \\
\texttt{ARGUMENT: \{'Dieser Vorschlag beinhaltet alle meine höchsten Werte (A48:868, A04:822, B21:808, B53:780, A45:668) und bleibt unter dem Aufwandlimit (3616). Es ist der optimale Deal für mich.'\}} \\
            }
        }
    }
    & \\ \\

    \theutterance \stepcounter{utterance}  
    & & & \multicolumn{4}{p{0.6\linewidth}}{
        \cellcolor[rgb]{0.9,0.9,0.9}{
            \makecell[{{p{\linewidth}}}]{
                \texttt{\tiny{[GM$\langle$P2]}}
                \texttt{STRATEGISCHE ÜBERLEGUNGEN: \{"Die letzte Vorschlag des Gegners enthält A48, das für mich nur 65 Wichtigkeit hat, und andere Items wie B21 und B53 mit niedriger Priorität für mich. Meine Gesamtwichtigkeit beträgt 3617, was unwahrscheinlich ist, dass weiterer Fortschritt möglich ist. Da ich keine bessere Vereinbarung erreichen kann, akzeptiere ich den Vorschlag, um das Spiel zu beenden."\}  } \\
\texttt{ZUSTIMMUNG: \{"B85", "C49", "A04", "A45", "A89", "B53", "B21", "A48", "C04"\}} \\
            }
        }
    }
     \\ \\

    \theutterance \stepcounter{utterance}  
    & & & \multicolumn{2}{p{0.3\linewidth}}{
        \cellcolor[rgb]{0.9,0.9,0.9}{
            \makecell[{{p{\linewidth}}}]{
                \texttt{\tiny{[GM$|$GM]}}
                \texttt{abort} \\
            }
        }
    }
    & & \\ \\

    \theutterance \stepcounter{utterance}  
    & & & \multicolumn{2}{p{0.3\linewidth}}{
        \cellcolor[rgb]{0.9,0.9,0.9}{
            \makecell[{{p{\linewidth}}}]{
                \texttt{\tiny{[GM$|$GM]}}
                \texttt{end game} \\
            }
        }
    }
    & & \\ \\

\end{supertabular}
}

\end{document}
